\begin{lemma}
Let \(\varepsilon>0\), and set
\[
\eta=\min\{\varepsilon/2,1/32\},
\]
\[
\varepsilon_1=
\frac12\min\left\{
\eta,\left(\frac{\eta^3}{32}\right)^2,\frac{\eta^3}{40}
\right\},
\qquad
\varepsilon_2=\frac12\min\left\{\varepsilon_1,\frac{\varepsilon_1}{30},1\right\}.
\]
Then there is \(n_0\) such that
\[
\varepsilon_2^{-2}\le n_0
\]
and
\[
\noderef{ParameterHierarchyEventualComparisons}(\eta,\varepsilon_1,
\varepsilon_2,n_0)
\]
holds.
\end{lemma}
\begin{proof}
By \(\noderef{ParameterHierarchyConstantBounds}\), the constants satisfy
\[
0<\eta<1/16,\qquad 0<\varepsilon_2<\varepsilon_1,\qquad
3\varepsilon_2\le\varepsilon_1/10,\qquad 0\le\varepsilon_2\le1.
\]
Choose \(n_0\) at least \(\varepsilon_2^{-2}\), and then enlarge it finitely
many times as described below.  We prove that every \(n>n_0\) satisfies all
conjuncts in \(\noderef{ParameterHierarchyEventualComparisons}\).

Write
\[
L=\log n,\qquad
k_R=(1+\eta)\sqrt{nL},\qquad
K=\noderef{TwoBiteNaturalIndependenceScale}(1+\eta,n),\qquad
k=(K:\mathbb R),
\]
\[
m_R=\frac{n}{L^2},\qquad
M=\noderef{TwoBiteNaturalReducedVertexCount}(n),\qquad
m=(M:\mathbb R),
\]
and
\[
p=\frac12\sqrt{\frac Ln},\qquad
t_1=\frac{\sqrt{nL}}{\log L},\qquad
t_2=n^{1/4+\eta},\qquad
t_3=n^{2\eta}.
\]
First enlarge \(n_0\) so that for every \(n>n_0\),
\[
n>1,\quad L>1,\quad \log L>0,\quad k_R>2,\quad m_R>2,\quad
\sqrt{L/n}\le1,\quad \frac{k_R+1}{n}\le1,\quad
\frac{\sqrt n}{2L^{3/2}}\ge1.
\]
These are eventual conditions because \(L\to\infty\), \(\log L\to\infty\),
\(\sqrt{nL}\to\infty\), \(n/L^2\to\infty\), \(\sqrt{L/n}\to0\),
\((k_R+1)/n\to0\), and \(\sqrt n/(2L^{3/2})\to\infty\).

For such \(n\), \(L>1\) makes \(m_R\) and \(k_R\) nonnegative.  Applying
\(\noderef{NaturalParameterApproximation}\) with \(\kappa=1+\eta\) gives
\[
k_R\le k<k_R+1,\qquad m\le m_R<m+1.
\]
Since \(k_R>2\), the ceiling inequality gives \(k<k_R+1<2k_R\), while
\((k_R+1)/n\le1\) gives \(k<n\) and hence \(K\le n\).  Since \(m_R>2\), the
floor inequality gives
\[
\frac12m_R\le m\le m_R,\qquad 0<m.
\]
Also \(p\ge0\) and \(p\le1/2\) follow from \(\sqrt{L/n}\ge0\) and
\(\sqrt{L/n}\le1\).  Finally,
\[
2pm\ge p m_R
=\frac12\sqrt{\frac Ln}\frac{n}{L^2}
=\frac{\sqrt n}{2L^{3/2}}\ge1.
\]
This proves the initial positivity, rounding, and \(1\le2pm\) conjuncts.

It remains to satisfy the finitely many asymptotic inequalities.  Each is
obtained by enlarging \(n_0\) so that a displayed ratio is below a fixed
positive constant or above a fixed positive constant.  The needed eventual
facts are all consequences of the standard growth rule
\[
\frac{(\log n)^a(\log\log n)^b}{n^c}\to0\quad(c>0),
\]
together with the reciprocal divergence statement.  We record the concrete
ratio for each conjunct.

For \(n^{4\eta}k\le(\varepsilon_1/2)k^2\), divide by \(k^2>0\) and use
\(k\ge k_R\).  It is enough that
\[
\frac{n^{4\eta}}{k_R}
=\frac{1}{1+\eta}\frac{n^{-1/2+4\eta}}{\sqrt L}
\le\frac{\varepsilon_1}{2},
\]
which is eventual because \(\eta<1/16\).  For
\[
2kL^2\le(\varepsilon_1/8)k^2,
\]
it is enough that
\[
\frac{2L^2}{k_R}
=\frac{2}{1+\eta}\frac{L^{3/2}}{\sqrt n}
\le\frac{\varepsilon_1}{8}.
\]

For the three-term \(\noderef{RealChooseTwo}\) comparison, use \(L>1\) and
\[
\frac1{\sqrt L}-\frac1{2L}\ge\frac1{2\sqrt L}.
\]
Since \(\noderef{RealChooseTwo}(k)\ge k^2/4\) for \(k>2\), the desired
inequality follows once
\[
\frac{16L^{5/2}}{k_R}\le1.
\]

For the exponential binomial inequality, \(K\le n\) gives
\[
\binom nK\le n^K=\exp(kL).
\]
Also \(\noderef{RealChooseTwo}(k)\ge k^2/4\), \(k\ge k_R\), and
\(m\le m_R=n/L^2\), so the exponent magnitude is at least
\[
\frac{k_R^4}{16Lm_R n^{8\eta}}
=\frac{(1+\eta)^4}{16}n^{1-8\eta}L^3.
\]
Enlarge \(n_0\) so this is at least \(2kL+L\).  Then
\[
\exp(-(2kL+L))\exp(kL)\le\exp(-L)=n^{-1}.
\]

We now verify the remaining conjuncts one by one, each time enlarging \(n_0\)
if necessary.  We use throughout
\[
k_R\le k<2k_R,\qquad \frac12m_R\le m\le m_R,\qquad
k_R=(1+\eta)\sqrt{nL},\qquad m_R=\frac n{L^2}.
\tag{A}
\]
For \(k>2\), the definition
\[
\noderef{RealChooseTwo}(k)=\frac{k(k-1)}2
\]
gives
\[
\frac{k^2}{4}\le \noderef{RealChooseTwo}(k)\le\frac{k^2}{2}. \tag{C}
\]
More generally, whenever \(x\ge0\), we will use the elementary upper bound
\[
\noderef{RealChooseTwo}(x)\le x^2. \tag{U}
\]
Finally, since \(k<2k_R\) and \(t_1=\sqrt{nL}/\log L\), after enlarging the
threshold so that \(\log L\ge1\),
\[
\frac{2k}{t_1}+1
<4(1+\eta)\log L+1
\le 5(1+\eta)\log L. \tag{D}
\]

For the conjunct
\[
\frac{2}{\sqrt L}\noderef{RealChooseTwo}(k)
\le \frac{\varepsilon_1}{2}k^2,
\]
the upper bound in (C) gives
\[
\frac{2}{\sqrt L}\noderef{RealChooseTwo}(k)
\le \frac{k^2}{\sqrt L}.
\]
Choose the threshold so that
\[
\frac1{\sqrt L}\le \frac{\varepsilon_1}{2}. \tag{T5}
\]
Then the desired inequality follows.

For the logarithmic conjunct
\[
4\log k\le \frac{\eta^3}{2}pk^2,
\]
use \(k<2k_R=2(1+\eta)\sqrt{nL}\), \(k\ge k_R\), and
\[
p k_R^2
=\frac12\sqrt{\frac Ln}\,(1+\eta)^2nL
=\frac{(1+\eta)^2}{2}\sqrt n\,L^{3/2}.
\]
Thus it is enough to require
\[
\frac{4\log(2(1+\eta)\sqrt{nL})}{p k_R^2}
=
\frac{8\log(2(1+\eta)\sqrt{nL})}
{(1+\eta)^2\sqrt n\,L^{3/2}}
\le \frac{\eta^3}{2}. \tag{T6}
\]
The left side tends to \(0\), so we enlarge \(n_0\) to make (T6) hold.

For
\[
t_1k\le \varepsilon_1 k^2,
\]
divide by the positive number \(k\).  Since \(k\ge k_R\),
\[
\frac{t_1}{k}\le\frac{t_1}{k_R}
=\frac{\sqrt{nL}/\log L}{(1+\eta)\sqrt{nL}}
=\frac1{(1+\eta)\log L}.
\]
Choose the threshold so that
\[
\frac1{(1+\eta)\log L}\le\varepsilon_1. \tag{T7}
\]

For
\[
\noderef{RealChooseTwo}\left(\frac{2k}{t_1}+1\right)\,100L^3
\le \frac12 k,
\]
combine (D) and (U):
\[
\noderef{RealChooseTwo}\left(\frac{2k}{t_1}+1\right)100L^3
\le 100\bigl(5(1+\eta)\log L\bigr)^2L^3.
\]
Since \(k\ge k_R\), it is enough to require
\[
\frac{200\bigl(5(1+\eta)\bigr)^2(\log L)^2L^3}
{(1+\eta)\sqrt{nL}}\le1. \tag{T8}
\]
Equivalently, the displayed ratio is a constant multiple of
\((\log L)^2L^{5/2}/\sqrt n\), so it tends to \(0\).

For
\[
\frac{2k}{t_1}(2pm)\le \frac{\varepsilon_1}{10}k,
\]
divide by \(k>0\).  Using \(m\le m_R\),
\[
2pm\le 2p m_R
=\sqrt{\frac Ln}\frac n{L^2}
=\frac{\sqrt n}{L^{3/2}}.
\]
Therefore
\[
\frac{1}{k}\frac{2k}{t_1}(2pm)
=\frac{2(2pm)}{t_1}
\le
\frac{2\sqrt n/L^{3/2}}{\sqrt{nL}/\log L}
=\frac{2\log L}{L^2}.
\]
Choose the threshold so that
\[
\frac{2\log L}{L^2}\le\frac{\varepsilon_1}{10}. \tag{T9}
\]

For
\[
\frac{2k}{t_1}\noderef{RealChooseTwo}(2pm)\le\varepsilon_1 k^2,
\]
use (U), \(2pm\le\sqrt n/L^{3/2}\), and \(k\ge k_R\).  After division by
\(k^2\), the left side is at most
\[
\frac{2}{t_1k_R}\left(\frac{\sqrt n}{L^{3/2}}\right)^2
=
\frac{2\log L}{(1+\eta)L^4}.
\]
Choose the threshold so that
\[
\frac{2\log L}{(1+\eta)L^4}\le\varepsilon_1. \tag{T10}
\]

For
\[
\noderef{RealChooseTwo}\left(\frac{2k}{t_1}+1\right)\,200L^3
\le \frac{\varepsilon_1}{10}k,
\]
the same estimate as in (T8) gives the sufficient condition
\[
\frac{2000\bigl(5(1+\eta)\bigr)^2(\log L)^2L^3}
{\varepsilon_1(1+\eta)\sqrt{nL}}\le1. \tag{T11}
\]

For
\[
\noderef{RealChooseTwo}\left(\frac{2k}{t_1}+1\right)
\noderef{RealChooseTwo}(200L^3)
\le \frac{\varepsilon_1}{10}k^2,
\]
use (D), (U), and
\(\noderef{RealChooseTwo}(200L^3)\le(200L^3)^2\).  Since
\(k^2\ge k_R^2=(1+\eta)^2nL\), it is enough that
\[
\frac{400000\bigl(5(1+\eta)\bigr)^2(\log L)^2L^5}
{\varepsilon_1(1+\eta)^2 n}\le1. \tag{T12}
\]

The next three conjuncts,
\[
3\varepsilon_2\le\varepsilon_1/10,\qquad
0\le\varepsilon_2,\qquad \varepsilon_2\le1,
\]
are exactly among the constant bounds supplied by
\(\noderef{ParameterHierarchyConstantBounds}\), so no threshold is needed.

For
\[
2pm\le\frac{\varepsilon_2}{100}t_1,
\]
the estimate \(2pm\le\sqrt n/L^{3/2}\) gives
\[
\frac{2pm}{t_1}\le
\frac{\sqrt n/L^{3/2}}{\sqrt{nL}/\log L}
=\frac{\log L}{L^2}.
\]
Choose the threshold so that
\[
\frac{\log L}{L^2}\le\frac{\varepsilon_2}{100}. \tag{T16}
\]

For
\[
\left(\frac{2k}{t_1}+1\right)100L^3
\le\frac{\varepsilon_2}{100}t_1,
\]
use (D).  It is enough that
\[
\frac{100\cdot5(1+\eta)(\log L)^2L^3}{\sqrt{nL}}
\le\frac{\varepsilon_2}{100}. \tag{T17}
\]
For
\[
\left(\frac{2k}{t_1}+1\right)200L^3
\le\frac{\varepsilon_2}{100}t_1,
\]
it is enough that
\[
\frac{200\cdot5(1+\eta)(\log L)^2L^3}{\sqrt{nL}}
\le\frac{\varepsilon_2}{100}. \tag{T18}
\]
For
\[
\left(\frac{2k}{t_1}+1\right)400L^5
\le\frac{\varepsilon_2}{100}t_1,
\]
it is enough that
\[
\frac{400\cdot5(1+\eta)(\log L)^2L^5}{\sqrt{nL}}
\le\frac{\varepsilon_2}{100}. \tag{T19}
\]
The three left sides in (T17)--(T19) tend to \(0\).

For
\[
t_2 n^{1/4-\eta}\le\varepsilon_1 k,
\]
unfold \(t_2=n^{1/4+\eta}\).  The left side is \(n^{1/2}\).  Since
\(k\ge k_R=(1+\eta)\sqrt n\sqrt L\), it is enough that
\[
\frac{1}{\varepsilon_1(1+\eta)\sqrt L}\le1. \tag{T20}
\]

For the cutoff comparison
\[
\noderef{TwoBiteLargeCutoff}(\eta,n)<\noderef{TwoBiteHugeCutoff}(n),
\]
unfold the definitions:
\[
\noderef{TwoBiteLargeCutoff}(\eta,n)=n^{1/4+\eta},
\qquad
\noderef{TwoBiteHugeCutoff}(n)=\frac{\sqrt{nL}}{\log L}=t_1.
\]
It is enough to require
\[
\frac{n^{1/4+\eta}\log L}{\sqrt{nL}}
=\frac{\log L}{n^{1/4-\eta}\sqrt L}<1. \tag{T21}
\]
This tends to \(0\), since \(\eta<1/16<1/4\).

For the huge-cutoff upper bound
\[
(1+\varepsilon_2)L^2\frac{t_2}{L}
  +L(1+\varepsilon_2)pm
<\noderef{TwoBiteHugeCutoff}(n),
\]
use \(\varepsilon_2\le1\), so \(1+\varepsilon_2\le2\).  The first term is at
most \(2Lt_2\), and the second term is at most
\[
2Lpm\le 2L\cdot\frac12\frac{\sqrt n}{L^{3/2}}
=\frac{\sqrt n}{\sqrt L}.
\]
Since \(\noderef{TwoBiteHugeCutoff}(n)=t_1=\sqrt{nL}/\log L\), it is enough to
choose the threshold so that
\[
\frac{2Lt_2}{t_1}
=\frac{2(\log L)L^{1/2}}{n^{1/4-\eta}}<\frac12,
\qquad
\frac{\sqrt n/\sqrt L}{t_1}
=\frac{\log L}{L}<\frac12. \tag{T22}
\]

For the large-cutoff lower bound
\[
100(n^{1/4}+1/2)L^3\le\noderef{TwoBiteLargeCutoff}(\eta,n),
\]
first enlarge the threshold so that \(n^{1/4}\ge1\).  Then
\[
n^{1/4}+1/2\le2n^{1/4},
\]
and it is enough that
\[
\frac{200L^3}{n^\eta}\le1. \tag{T23}
\]

For
\[
k< n^{1/4}t_2
-\noderef{RealChooseTwo}(n^{1/4})\,100L^3,
\]
we use (U) to get
\[
\noderef{RealChooseTwo}(n^{1/4})\,100L^3\le100n^{1/2}L^3.
\]
Choose the threshold so that
\[
\frac{100L^3}{n^\eta}\le\frac12. \tag{T24a}
\]
Then the negative term is at most \(\frac12n^{1/2+\eta}\), so the right side
is at least \(\frac12n^{1/2+\eta}\).  Also \(k<2k_R\), so it is enough to
require
\[
\frac{2(1+\eta)\sqrt L}{n^\eta}<\frac12. \tag{T24b}
\]

For
\[
t_3^2k\le\frac{\varepsilon_1}{2}k^2,
\]
unfold \(t_3=n^{2\eta}\), so \(t_3^2=n^{4\eta}\).  This is the same reduction
as the first \(n^{4\eta}\)-conjunct: it follows from
\[
\frac{n^{4\eta}}{k_R}\le\frac{\varepsilon_1}{2}. \tag{T25}
\]

For
\[
k\left(\frac{2k}{L}+n^{1/4}\right)\le\varepsilon_1 k^2,
\]
divide by \(k>0\).  It is enough to put each summand below
\((\varepsilon_1/2)k\).  The first is controlled by
\[
\frac{2}{L}\le\frac{\varepsilon_1}{2}, \tag{T26a}
\]
and the second by \(k\ge k_R\):
\[
\frac{n^{1/4}}{k_R}
=\frac{1}{(1+\eta)n^{1/4}\sqrt L}
\le\frac{\varepsilon_1}{2}. \tag{T26b}
\]

For the combined coefficient inequality
\[
8\sqrt{\varepsilon_1}pk^2+10\varepsilon_1pk^2+4\log k
\le \eta^3pk^2,
\]
the constant bound gives
\[
8\sqrt{\varepsilon_1}+10\varepsilon_1\le\eta^3/2.
\]
Multiplying by the nonnegative quantity \(pk^2\) gives the first two terms at
most \((\eta^3/2)pk^2\).  The logarithmic conjunct already proved from (T6)
gives
\[
4\log k\le(\eta^3/2)pk^2.
\]
Adding the two estimates gives the desired \(\eta^3pk^2\) bound.

For the first final exponential,
\[
\exp\!\left(kL+2kn^{1/4-3\eta}L-n^{3/4-2\eta}\right)\le\varepsilon_1,
\]
use \(k<2k_R=2(1+\eta)\sqrt n\sqrt L\).  The positive part of the exponent is
bounded by
\[
2(1+\eta)n^{1/2}L^{3/2}
+4(1+\eta)n^{3/4-3\eta}L^{3/2}.
\]
Choose the threshold so that
\[
2(1+\eta)n^{1/2}L^{3/2}\le\frac14 n^{3/4-2\eta},
\qquad
4(1+\eta)n^{3/4-3\eta}L^{3/2}\le\frac14 n^{3/4-2\eta}. \tag{T28a}
\]
These are equivalent to requiring fixed multiples of
\(L^{3/2}/n^{1/4-2\eta}\) and \(L^{3/2}/n^\eta\) to be at most \(1/4\), and
both ratios tend to \(0\).  Under (T28a), the exponent is at most
\[
-\frac12 n^{3/4-2\eta}.
\]
Since \(\varepsilon_1>0\), choose the threshold further so that
\[
\exp\!\left(-\frac12 n^{3/4-2\eta}\right)\le\varepsilon_1. \tag{T28b}
\]

For the last exponential,
\[
\exp\!\left(-\frac{\eta}{4}kL\right)\le\varepsilon_1,
\]
use \(k\ge k_R\) to get
\[
\frac{\eta}{4}kL
\ge \frac{\eta(1+\eta)}4\sqrt n\,L^{3/2}.
\]
The right side tends to infinity, so enlarge the threshold so that
\[
\exp\!\left(-\frac{\eta(1+\eta)}4\sqrt n\,L^{3/2}\right)\le\varepsilon_1.
\tag{T29}
\]

All displayed threshold requirements (T5)--(T29) are finite: each ratio that
must be small is a fixed constant times
\((\log n)^a(\log\log n)^b/n^c\) with \(c>0\), and each exponent magnitude
that must be large is a positive power of \(n\) divided by a fixed power of
\(\log n\) and \(\log\log n\).  Enlarge \(n_0\) to dominate all of these
finitely many thresholds.

With this single enlarged \(n_0\), every conjunct in
\(\noderef{ParameterHierarchyEventualComparisons}(\eta,\varepsilon_1,
\varepsilon_2,n_0)\) holds for all \(n>n_0\), and by construction
\(\varepsilon_2^{-2}\le n_0\).
\end{proof}
